\section{Threat model and system model}
\label{sec:threat}

A security paper states its threat model before it taxonomizes. This paper
spans two literatures that assume different ones, so we state both and say
where each applies.

\subsection{System model}

The subject is an LLM agent: a model that, within one task, repeatedly reads
content it did not author, decides, and acts. Reading happens through the nine
channels of Section~\ref{sec:taxonomy}. Acting means emitting text to a user or
calling a tool whose effects leave the system.

Three properties separate this from a single-turn chatbot, and the taxonomy
keys on all three. Content from many sources shares one context window, with
no trustworthy marking of where any of it came from. The agent's own output
can become another agent's input, or its own input a turn later through
memory. And an action taken on corrupted context cannot be undone by fixing
the context afterwards.

\subsection{Adversarial threat model}

The adversary controls the content of at least one channel. They do not
control the model weights, the decoding procedure, or the agent's code. This
is the standing assumption of the security track, and it covers the realistic
deployment cases. An adversary can write a document a retriever may fetch,
publish a tool whose description the agent reads, post content a browsing
agent visits, or contribute a skill file the agent loads.

The adversary knows the defense and may adapt the payload to it. That is
precisely why average-case evaluation does not transfer. The adversary does
not need to know the user's query: every attack we run is query-agnostic by
construction.

We exclude three things. Training-time poisoning and weight-level attacks do
not travel through the context window. A user attacking their own agent
involves no third-party content, so the trust boundary we study does not
exist. Denial of service through resource exhaustion sits in the taxonomy as
\meth{resource-abuse}, has no entries, and we make no claims about it.

\subsection{Incidental threat model}

There is no adversary. Content goes wrong through ordinary pipeline behavior.
A retrieved document is stale. A passage is irrelevant but topically close.
Important text sits in the middle of a long input. A retrieved fact
contradicts what the model learned in training. The machine learning track
assumes this model and evaluates on average-case benchmarks, which is the
right choice when nothing is adapting.

\subsection{Why both, and what connects them}

The two models differ in exactly one assumption: adaptivity. They share the
channel, the position in the context window, and the downstream consequence.
Section~\ref{sec:taxonomy} shows that 97\% of published defenses are validated
against one model and not the other.

\textbf{This paper asks whether that single validation transfers.} The
assumption behind it is reasonable. A method that makes a model robust to a
stale document is working on the same surface an attacker writes to. It has
simply never been tested at scale. Section~\ref{sec:experiments} tests it.
